\chapter{Transition vers l'implémentation retenue}
\label{chap:transition-implementation}
\justifying
\sloppy

\section{De la structure de référence à une architecture exploitable}

L'architecture Vernier multiphase de référence constitue une base théorique et architecturale solide pour obtenir une résolution temporelle fine \cite{annagrebah2021thesis}. Elle ne peut toutefois pas être reprise telle quelle dans une chaîne de lecture numérique complète. Dans \textbf{SPADMIC}, le problème ne consiste pas seulement à démontrer le potentiel d'un \textbf{TDC} Vernier multiphase : il faut organiser la capture des événements, le stockage des résultats, leur transmission et leur exploitation hors ASIC dans un prototype complet.

La continuité avec l'architecture Vernier multiphase de référence se situe au niveau des principes fondamentaux :
\begin{itemize}
    \item utilisation de deux oscillateurs proches pour définir une résolution issue d'une différence de délai ;
    \item séparation entre information grossière et information fine ;
    \item observation multiphase par matrice de détection ;
    \item nécessité d'une calibration adaptée pour transformer les observables brutes en temps exploitable.
\end{itemize}

La différence majeure vient du contrat d'intégration. Le bloc \textbf{MPTDC} n'est pas une cellule isolée : l'implémentation active est construite autour d'une matrice de détection \mbox{$8\times8$}, d'une capture asynchrone, d'une banque de contextes et d'une chaîne de lecture dans \texttt{clk\_sys}. Ces choix préservent le principe Vernier tout en imposant une architecture routable, réarmable et compatible avec un prototype en \SI{180}{\nano\meter}.

\section{Évolutions nécessaires dans le contexte SPADMIC}

L'architecture développée s'écarte d'une simple preuve de concept sur plusieurs aspects structurants. Elle doit répondre simultanément à des besoins de robustesse, d'intégration numérique et de lecture :
\begin{itemize}
    \item prise en charge d'événements \texttt{START}/\texttt{STOP} avec une logique asynchrone explicite ;
    \item séparation entre la capture rapide et la lecture système ;
    \item format de sortie numérique compatible avec une chaîne de lecture externe ;
    \item calibration pensée comme une opération externe à l'ASIC, et non comme une correction embarquée figée dans le silicium.
\end{itemize}

Ces choix déplacent le centre de gravité du problème. Le principe Vernier fournit le mécanisme de résolution fine ; l'implémentation active doit ensuite garantir que cette information reste cohérente lorsqu'elle traverse des domaines temporels différents, qu'elle est figée dans un contexte, puis qu'elle est sérialisée vers une sortie compacte.

\section{Vision globale de l'architecture active}

Avant d'isoler les arbitrages de conception, il faut placer les blocs dans leur topologie numérique globale. La figure \ref{fig:vernier-transition-map} reprend le formalisme matériel de l'architecture de référence et montre les transformations essentielles : la prélogique devient un front-end de capture asynchrone, la matrice de phase est ramenée à \mbox{$8\times8$}, la mémoire d'échantillonnage devient une banque de contextes, et une chaîne de lecture système est ajoutée en sortie.

\begin{figure}[H]
    \centering
    \includegraphics[width=\textwidth]{tdc/vernier_transition_map.pdf}
    \caption{Transformation de l'architecture Vernier de référence vers l'implémentation active \textbf{MPTDC}.}
    \label{fig:vernier-transition-map}
\end{figure}

Cette transition doit aussi être lue sous l'angle des domaines temporels. L'architecture réelle superpose une zone d'événements asynchrones, deux familles de signaux issues des oscillateurs lent et rapide, puis un domaine système \texttt{clk\_sys} chargé du séquençage de capture, de la lecture du contexte et de la sérialisation. La figure \ref{fig:clock-domain-diagram} fournit cette carte avant l'analyse des compromis.

\begin{figure}[H]
    \centering
    \makebox[\textwidth][c]{%
        \includegraphics[width=1.05\textwidth]{tdc/clock_domain_diagram.pdf}%
    }
    \caption{Carte des domaines temporels et des traversées CDC de l'architecture MPTDC active.}
    \label{fig:clock-domain-diagram}
\end{figure}

Un \textbf{TDC} Vernier actif ne possède pas un seul domaine synchrone. Les fronts \texttt{START} et \texttt{STOP} arrivent sans relation de phase garantie avec \texttt{clk\_sys}. Les anneaux lent et rapide créent ensuite leurs propres domaines temporels, directement liés aux oscillateurs. Le domaine \texttt{clk\_sys}, beaucoup plus lent, sert à séquencer la fermeture de mesure, figer l'image utile, écrire le contexte et lire la trame de sortie. Cette structure interdit de traiter la conception comme une FSM unique entourée de registres synchrones ordinaires.

Le transport direct d'un bus multi-bit entre domaines serait dangereux. Chaque bit peut être routé différemment, subir un délai de propagation distinct et atteindre le registre de destination à un instant légèrement différent. Au moment de l'échantillonnage, certains bits peuvent donc appartenir à l'ancienne valeur et d'autres à la nouvelle. La valeur obtenue peut être incohérente même si chaque bit traverse individuellement une chaîne de synchronisation. Ce risque est particulièrement critique pour un compteur binaire, car plusieurs bits peuvent changer sur une même transition.

Pour les compteurs grossiers, le choix retenu consiste à transporter des codes Gray. Un code Gray limite la transition logique à un seul bit à la fois ; l'ambiguïté résiduelle est alors bornée à la valeur précédente ou suivante, au lieu de produire une combinaison binaire impossible. Les phases rapides et lentes ne sont pas transportées comme des bus synchrones continus vers \texttt{clk\_sys}. L'architecture transporte plutôt des commandes compactes, des compteurs Gray, des snapshots, des images figées et des impulsions synchronisées.

La figure~\ref{fig:gray-cdc-principle} illustre pourquoi le transfert direct d'un compteur multi-bit entre domaines d'horloge est dangereux, puis comment le codage Gray et la synchronisation par deux flip-flops limitent ce risque.

\begin{figure}[H]
    \centering
    \makebox[\textwidth][c]{\resizebox{1.00\textwidth}{!}{%
    \begin{tikzpicture}[x=1cm,y=1cm,>=Latex]
        \tikzset{
            panel/.style={draw=black!28,rounded corners=2pt,fill=black!2,line width=0.85pt},
            title/.style={font=\sffamily\bfseries\normalsize},
            bitlabel/.style={font=\ttfamily\small,anchor=east},
            wave/.style={line width=1.25pt,draw=black},
            sample/.style={draw=black!70,densely dashed,line width=1.05pt},
            note/.style={font=\sffamily\footnotesize,align=center},
            noteBox/.style={font=\sffamily\scriptsize,align=center,fill=white,inner sep=1.6pt,rounded corners=1pt},
            warn/.style={noteBox,text=red!85!black},
            oknote/.style={noteBox,text=blue!75!black},
            good/.style={noteBox,text=green!45!black},
            ff/.style={draw=blue!70!black,fill=blue!7,rounded corners=1pt,
                       minimum width=1.16cm,minimum height=0.78cm,font=\sffamily\small\bfseries},
            arr/.style={-{Latex[length=2.1mm]},line width=1.15pt,draw=blue!75!black},
            clkline/.style={draw=black!65,line width=0.95pt}
        }

        \node[panel,minimum width=4.90cm,minimum height=4.65cm,anchor=south west] at (0,0) {};
        \node[panel,minimum width=4.90cm,minimum height=4.65cm,anchor=south west] at (5.45,0) {};
        \node[panel,minimum width=4.90cm,minimum height=4.65cm,anchor=south west] at (10.90,0) {};

        \begin{scope}[shift={(0,0)}]
            \node[title] at (2.45,4.32) {Compteur binaire};
            \node[font=\ttfamily\Large] at (1.24,3.78) {011};
            \draw[-{Latex[length=2.1mm]},red!85!black,line width=1.25pt] (1.76,3.78) -- (2.42,3.78);
            \node[font=\ttfamily\Large] at (2.95,3.78) {100};
            \node[warn] at (2.45,3.34) {plusieurs bits changent};

            \node[bitlabel] at (0.72,2.78) {b2};
            \node[bitlabel] at (0.72,2.18) {b1};
            \node[bitlabel] at (0.72,1.58) {b0};
            \draw[wave] (0.98,2.78) -- (1.60,2.78) -- (1.60,3.08) -- (4.35,3.08);
            \draw[wave] (0.98,2.48) -- (1.94,2.48) -- (1.94,2.18) -- (4.35,2.18);
            \draw[wave] (0.98,1.88) -- (2.28,1.88) -- (2.28,1.58) -- (4.35,1.58);
            \foreach \x/\yA/\yB in {1.60/2.78/3.08,1.94/2.48/2.18,2.28/1.88/1.58} {
                \draw[red!85!black,line width=1.45pt] (\x,\yA) -- (\x,\yB);
            }
            \draw[sample] (2.10,1.28) -- (2.10,3.16);
            \node[note,anchor=north,text=black!65,fill=white,inner sep=1.2pt] at (2.10,1.18) {front\\\texttt{clk\_sys}};
            \node[draw=red!65!black,fill=red!5,rounded corners=1pt,
                  font=\ttfamily\small,minimum width=0.90cm,minimum height=0.44cm] at (3.55,2.18) {001?};
            \node[warn] at (2.45,0.34) {risque : mot incoh\'erent};
        \end{scope}

        \begin{scope}[shift={(5.45,0)}]
            \node[title] at (2.45,4.32) {Compteur Gray};
            \node[font=\ttfamily\Large] at (1.24,3.78) {010};
            \draw[-{Latex[length=2.1mm]},blue!75!black,line width=1.25pt] (1.76,3.78) -- (2.42,3.78);
            \node[font=\ttfamily\Large] at (2.95,3.78) {110};
            \node[oknote] at (2.45,3.34) {un seul bit change};

            \node[bitlabel] at (0.72,2.78) {g2};
            \node[bitlabel] at (0.72,2.18) {g1};
            \node[bitlabel] at (0.72,1.58) {g0};
            \draw[wave] (0.98,2.78) -- (1.88,2.78) -- (1.88,3.08) -- (4.35,3.08);
            \draw[wave] (0.98,2.48) -- (4.35,2.48);
            \draw[wave] (0.98,1.58) -- (4.35,1.58);
            \draw[blue!75!black,line width=1.45pt] (1.88,2.78) -- (1.88,3.08);
            \draw[sample] (2.10,1.28) -- (2.10,3.16);
            \node[note,anchor=north,text=black!65,fill=white,inner sep=1.2pt] at (2.10,1.18) {front\\\texttt{clk\_sys}};

            \node[good] at (2.45,0.34) {ancienne ou nouvelle valeur};
        \end{scope}

        \begin{scope}[shift={(10.90,0)}]
            \node[title] at (2.45,4.32) {Synchroniseur 2 flip-flops};
            \node[font=\sffamily\scriptsize,anchor=east,text=black!70] at (1.02,3.20) {bit Gray};
            \node[font=\sffamily\scriptsize,anchor=east,text=black!70] at (1.02,2.98) {asynchrone};
            \node[ff] (ff1) at (1.70,2.62) {DFF1};
            \node[ff] (ff2) at (3.08,2.62) {DFF2};
            \node[font=\sffamily\tiny,text=blue!70!black] at ([xshift=-0.38cm,yshift=0.25cm]ff1.center) {D};
            \node[font=\sffamily\tiny,text=blue!70!black] at ([xshift=0.38cm,yshift=0.25cm]ff1.center) {Q};
            \node[font=\sffamily\tiny,text=blue!70!black] at ([xshift=-0.38cm,yshift=0.25cm]ff2.center) {D};
            \node[font=\sffamily\tiny,text=blue!70!black] at ([xshift=0.38cm,yshift=0.25cm]ff2.center) {Q};

            \draw[arr] (0.72,2.62) -- (ff1.west);
            \draw[arr] (ff1.east) -- (ff2.west);
            \draw[arr] (ff2.east) -- (4.48,2.62);
            \node[font=\sffamily\scriptsize,anchor=west,text=black!70] at (3.82,2.94) {bit stable};

            \draw[-{Latex[length=2.0mm]},clkline]
                (0.76,1.45) -- node[below,font=\sffamily\scriptsize] {\texttt{clk\_sys}} (4.18,1.45);
            \draw[clkline] (1.70,1.45) -- (1.70,2.20);
            \draw[clkline] (3.08,1.45) -- (3.08,2.20);
            \draw[clkline] (1.56,2.20) -- (1.70,2.35) -- (1.84,2.20);
            \draw[clkline] (2.94,2.20) -- (3.08,2.35) -- (3.22,2.20);

            \node[oknote] at (2.45,0.76) {appliqu\'e bit par bit};
            \node[good] at (2.45,0.34) {stabilisation avant usage};
        \end{scope}
    \end{tikzpicture}%
    }}
    \caption{Principe de traversée CDC d'un compteur : risque du codage binaire, intérêt du Gray et synchronisation par deux flip-flops.}
    \label{fig:gray-cdc-principle}
\end{figure}

La synchronisation à deux flip-flops consiste à faire traverser chaque bit par deux bascules cadencées par le domaine de destination. La première bascule peut devenir métastable si le bit change trop près du front d'horloge ; la seconde laisse un cycle supplémentaire pour stabiliser le signal avant son utilisation par la logique \(clk\_sys\). Ce mécanisme ne rend pas la traversée instantanée et ne permet pas de transporter arbitrairement un bus binaire multi-bit. Il devient robuste lorsqu'il est associé à un codage Gray, car une seule transition logique est possible entre deux valeurs successives.

La méthode RTL appliquée combine donc plusieurs mécanismes complémentaires. Les compteurs traversent la frontière sous forme Gray et sont resynchronisés dans le domaine de destination. Les fronts \texttt{START}/\texttt{STOP} déclenchent des verrous asynchrones intentionnels afin de ne pas injecter \texttt{clk\_sys} sur le chemin temporel mesuré. La fin de conversion ne lance pas immédiatement la lecture : elle déclenche d'abord un snapshot de l'image de mesure, puis une écriture dans un contexte. Le clear est ordonné après cette protection, ce qui évite de détruire la matrice source avant que l'image stable soit disponible pour la lecture.

\section{Choix d'intégration et compromis d'architecture}

Le passage du principe Vernier à un bloc ASIC intégrable impose de figer une architecture qui reste suffisamment riche pour la calibration, mais suffisamment compacte pour être routée et intégrée dans le prototype \textbf{SPADMIC}. L'implémentation active retient ainsi une matrice de phase \mbox{$8\times8$}, deux contextes de capture et une sortie numérique étroite. La matrice fournit les couples de phase nécessaires à l'observation fine, tout en limitant la surface, la charge appliquée aux sorties de phase des oscillateurs, la complexité de lecture du contexte et l'effort de fermeture temporelle.

La capture des événements reste volontairement asynchrone. Les signaux \texttt{START} et \texttt{STOP} ne sont pas d'abord ramenés dans \texttt{clk\_sys}, car cela déplacerait l'incertitude temporelle sur le chemin même que le \textbf{TDC} doit mesurer. Le front-end déclenche donc l'oscillateur lent au plus près de \texttt{START}, puis l'oscillateur rapide au plus près de \texttt{STOP}. Le domaine système intervient ensuite pour protéger l'image de mesure, orchestrer le nettoyage et préparer la lecture, sans transformer les phases analogiques en bus synchrones continus.

\begin{figure}[H]
    \centering
    \includegraphics[width=0.96\textwidth]{tdc/async_capture_timing.pdf}
    \caption{Capture asynchrone des événements \texttt{START}/\texttt{STOP} et protection de la fenêtre de mesure par \texttt{clk\_sys}.}
    \label{fig:async-capture-timing}
\end{figure}

Un point d'attention concerne la logique placée avant le déclenchement effectif des anneaux. Toute combinatoire sur le chemin \texttt{START}/\texttt{STOP} peut introduire un délai systématique et une contribution de jitter avant \texttt{slow\_en} et \texttt{fast\_en}. Cet effet reste peu visible sur une conversion isolée, mais il devient critique lors de mesures répétées ou d'une caractérisation fine. Il devra donc être évalué avec les blocs réels, puis en simulation post-layout, afin de quantifier la contribution ajoutée par le front-end de capture.

Le compromis le plus structurant concerne le découplage capture/lecture. Si la capture et la lecture restent couplées, le temps mort dépend directement du temps nécessaire à parcourir la matrice, construire les enregistrements et sérialiser la trame de sortie. Le choix RTL retenu consiste à figer la mesure par snapshot dans un contexte disponible, puis à nettoyer la matrice pour autoriser un nouvel événement. La lecture du contexte peut ensuite progresser dans \texttt{clk\_sys}, à un rythme compatible avec la FIFO et la sortie \SI{16}{\bit}.

\begin{figure}[H]
    \centering
    \includegraphics[width=0.95\textwidth]{tdc/context_pipeline.pdf}
    \caption{Pipeline à deux contextes permettant de recouvrir capture, lecture et sérialisation.}
    \label{fig:context-pipeline}
\end{figure}

Le double contexte ne signifie pas que la matrice physique de détecteurs de phase est dupliquée. Il correspond à deux zones de stockage d'images de mesure. Lors du snapshot, les informations utiles d'une conversion terminée, bitmap de matrice, compteurs, métadonnées et statut, sont figées dans un contexte disponible. Le gain PPA attendu se situe d'abord sur le temps mort : la capture n'attend plus la fin de la sérialisation. Le coût associé est une surface mémoire supplémentaire pour deux images de contexte, acceptée ici parce qu'elle reste bornée et qu'elle évite une file profonde plus coûteuse.

Le goulot est déplacé vers l'occupation des deux contextes lorsque la sortie aval ralentit. Si le drain ou le \textbf{TX} reste bloqué trop longtemps, les deux contextes peuvent devenir indisponibles et un nouveau \texttt{START} doit être rejeté de manière contrôlée. Le risque est maîtrisé par les états d'occupation, les statuts de contexte, la rétropression, le rejet explicite des nouvelles captures et le dimensionnement de la sortie. Le compromis retenu garde donc une surface de stockage limitée tout en rendant visible la saturation éventuelle.

La sortie adopte enfin un protocole compact sur \SI{16}{\bit}, sous forme d'une trame de mesure structurée. Le paquet contient un en-tête, des mots de données associés aux enregistrements de détection et un mot de fin. Il transporte les informations nécessaires à la reconstruction externe : valeurs de mesure, métadonnées de contexte, indicateurs de frontière et champs de statut. Le contrat de sortie reste unique et compact, sans introduire plusieurs chemins concurrents. La correction finale est déplacée hors ASIC : le silicium exporte une description stable et suffisamment riche de la mesure, tandis que la LUT et les modèles de calibration restent évolutifs côté hôte.

\begin{figure}[H]
    \centering
    \includegraphics[width=0.93\textwidth]{tdc/calibration_export_flow.pdf}
    \caption{Flux de sortie compact et calibration externe par décodage hôte et LUT.}
    \label{fig:calibration-export-flow}
\end{figure}

\begin{summaryBox}
La transition vers l'implémentation transforme le principe Vernier en un bloc ASIC complet : capture asynchrone, matrice de phase compacte, frontières CDC explicites, double contexte, lecture \texttt{clk\_sys}, paquet \SI{16}{\bit} et calibration externe forment un même compromis d'intégration.
\end{summaryBox}

\clearpage
